\section{再论闭区间上连续函数的性质}

前面讲述过闭区间上连续函数的性质，它们在微积分的理论中起着基本的作用．当时为了避免运用过多的概念，所有证明都是用区间套定理直接给出的．现在我们已经有了很多的工具，每个定理都可以给出更为简洁的证明．其实证明是很多的，读者可以尝试寻找其他证明，作为对自己的一种训练．

{\thmheadfont 有界性定理}\quad 若 $f(x)$ 在闭区间 $[a, b]$ 上连续，则 $f(x)$ 在 $[a, b]$ 有界．

{\thmheadfont 证明一}\quad 利用区间套定理，用反证法．如果不然，设 $f(x)$ 在 $[a, b]$ 无界，记 $a_1 = a$，$b_1 = b$．二等分 $[a, b]$，则在两个子区间中，至少有一个，$f(x)$ 在其上是无界的，记此子区间为 $[a_2, b_2]$．再二等分 $[a_2, b_2]$，$\cdots$，如此继续下去，便得一区间套 $\{[a_n, b_n]\}$，$f(x)$ 在其中的每个区间都是无界的．根据区间套定理，存在 $r \in [a, b]$，使得 $\lim\limits_{n \to \infty} a_n = \lim\limits_{n \to \infty} b_n = r$．由 $f(x)$ 在 $r$ 连续，知存在 $\delta > 0$，使得 $f(x)$ 在 $(r - \delta, r + \delta) \cap [a, b]$ 有界．但当 $n$ 充分大时，有 $[a_n, b_n] \subset (r - \delta, r + \delta) \cap [a, b]$，与 $f(x)$ 在 $[a_n, b_n]$ 无界矛盾，定理获证．

{\thmheadfont 证明二}\quad 利用紧致性定理，用反证法．如果不然，设 $f(x)$ 在 $[a, b]$ 无界，这时对任意正整数 $n$，存在 $x_n \in [a, b]$，使得 $|f(x_n)| \geqslant n$．由闭区间 $[a, b]$ 的紧致性，知存在 $\{x_n\}$ 的子数列 $\{x_{n_k}\}$，$\lim\limits_{k \to \infty} x_{n_k} = r \in [a, b]$，因此，$\lim\limits_{k \to \infty} f(x_{n_k}) = \infty$，与 $f(x)$ 在 $r$ 连续矛盾，定理获证．

{\thmheadfont 介值定理}\quad 若 $f(x)$ 在 $[a, b]$ 连续，$f(a) > 0$，$f(b) < 0$，则存在 $c \in [a, b]$，使得 $f(c) = 0$．

{\thmheadfont 证明}\quad 用有限覆盖定理．用反证法，若不然，任意 $x_0 \in [a, b]$，有 $f(x_0) \neq 0$，由 $f(x)$ 在 $x_0$ 点连续及保号性，存在 $\delta_{x_0} > 0$，当 $x \in (x_0 - \delta_{x_0}, x_0 + \delta_{x_0}) \cap [a, b]$ 时，$f(x)$ 不变号．令
\[
E = \{ (x_0 - \delta_{x_0}, x_0 + \delta_{x_0}) \mid x_0 \in [a, b] \},
\]
则 $E$ 是 $[a, b]$ 的一个覆盖，由有限覆盖定理，存在 $x_1, x_2, \cdots, x_n \in [a, b]$，使得
\[
[a, b] \subset \bigcup_{i=1}^{n} (x_i - \delta_{x_i}, x_i + \delta_{x_i}).
\]
不妨设 $x_1 < x_2 < \cdots < x_n$，且 $|i - j| \geqslant 2$ 时有
\[
(x_i - \delta_{x_i}, x_i + \delta_{x_i}) \cap (x_j - \delta_{x_j}, x_j + \delta_{x_j}) = \varnothing,
\]
于是任意 $i = 1, 2, \cdots, n-1$，有
\[
(x_i - \delta_{x_i}, x_i + \delta_{x_i}) \cap (x_{i+1} - \delta_{x_{i+1}}, x_{i+1} + \delta_{x_{i+1}}) \neq \varnothing.
\]
由于 $f(x)$ 在 $(x_i - \delta_{x_i}, x_i + \delta_{x_i})$ 不变号，则 $f(x)$ 在 $[a, b]$ 不变号，这与 $f(a) > 0$，$f(b) < 0$ 矛盾．故存在 $c \in [a, b]$，使得 $f(c) = 0$．

{\thmheadfont 最大值定理}\quad 若 $f(x)$ 在 $[a, b]$ 连续，则 $f(x)$ 在 $[a, b]$ 上达到其最大值与最小值．

{\thmheadfont 证明}\quad 只证最大值的情形．由于 $f(x)$ 在 $[a, b]$ 有界，则上确界
\[
M = \sup_{a \leqslant x \leqslant b} f(x)
\]
存在．由上确界性质，对任意 $n$，存在 $x_n \in [a, b]$，使得
\[
M \geqslant f(x_n) > M - \frac{1}{n}.
\]
由 $[a, b]$ 的紧致性知，存在子数列 $\{x_{n_k}\}$：$\lim\limits_{k \to \infty} x_{n_k} = r \in [a, b]$．在不等式
\[
M \geqslant f(x_{n_k}) > M - \frac{1}{n_k}
\]
中，令 $k \to \infty$ 取极限，便得 $f(r) = M$，定理证完．

{\thmheadfont 一致连续性定理}\quad 若 $f(x)$ 在 $[a, b]$ 连续，则 $f(x)$ 在 $[a, b]$ 一致连续．

{\thmheadfont 证明一}\quad 用有限覆盖定理．对任给 $\varepsilon > 0$，对任给 $x \in [a, b]$，存在 $\delta_x > 0$，使得只要 $x - \delta_x < t < x + \delta_x$，有
\[
|f(t) - f(x)| < \frac{\varepsilon}{2}.
\]
于是 $\left\{ \left. \left( x - \frac{\delta_x}{2}, x + \frac{\delta_x}{2} \right) \right| x \in [a, b] \right\}$ 构成了 $[a, b]$ 的一个覆盖．由有限覆盖定理，知存在有限个点 $x_1, \cdots, x_n$，使得
\[
\bigcup_{j=1}^{n} \left( x_j - \frac{\delta_j}{2}, x_j + \frac{\delta_j}{2} \right) \supset [a, b].
\]
其中 $\delta_j = \delta_{x_j}$．令 $\delta = \min\left( \dfrac{\delta_1}{2}, \cdots, \dfrac{\delta_n}{2} \right)$，则只要 $|x' - x''| < \delta$，总有 $x_k$，使 $x' \in \left( x_k - \dfrac{\delta_k}{2}, x_k + \dfrac{\delta_k}{2} \right)$，从而 $x'' \in (x_k - \delta_k, x_k + \delta_k)$，因此
\[
\begin{aligned}
|f(x') - f(x'')| &\leqslant |f(x') - f(x_k)| + |f(x_k) - f(x'')| \\
&< \frac{\varepsilon}{2} + \frac{\varepsilon}{2} = \varepsilon,
\end{aligned}
\]
这就证明了 $f(x)$ 在 $[a, b]$ 一致连续，定理证完．

{\thmheadfont 证明二}\quad 用反证法．如果不然，设 $f(x)$ 在 $[a, b]$ 上不一致连续，这时，存在 $\varepsilon_0 > 0$，对任意的 $\delta > 0$，存在 $x', x'' \in [a, b]$，使得 $|x' - x''| < \delta$，但
\[
|f(x') - f(x'')| \geqslant \varepsilon_0.
\]
分别取 $\delta = \dfrac{1}{n}$，$n = 1, 2, \cdots$，相应地存在 $x_n', x_n'' \in [a, b]$ 满足 $|x_n' - x_n''| < \dfrac{1}{n}$，但
\[
|f(x_n') - f(x_n'')| \geqslant \varepsilon.
\]
由 $[a, b]$ 的紧致性，存在 $\{x_n'\}$ 的子数列 $\{x_{n_k}'\}$，使 $\lim\limits_{k \to \infty} x_{n_k}' = \gamma \in [a, b]$．根据 $|x_{n_k}' - x_{n_k}''| < \dfrac{1}{n_k}$，知也有 $\lim\limits_{k \to \infty} x_{n_k}'' = \gamma$，从而 $\lim\limits_{k \to \infty} f(x_{n_k}') = \lim\limits_{k \to \infty} f(x_{n_k}'') = f(\gamma)$．在不等式
\[
|f(x_{n_k}') - f(x_{n_k}'')| \geqslant \varepsilon_0
\]
中令 $k \to \infty$ 取极限，便得 $0 \geqslant \varepsilon_0$，与 $\varepsilon_0 > 0$ 矛盾，故 $f(x)$ 在 $[a, b]$ 一致连续，定理证完．

上述这些性质，对开区间（或半开区间）的连续函数都是不一定成立的，其根源在于闭区间在实数系是紧致的（而且是连通的），而开区间或半开区间（虽然是连通的）却不是紧致的，它们和闭区间在实数系中虽只差一、两个点，但性质却迥然不同，真是“差之毫厘，谬之千里”．

还值得指出的是，介值定理与最值定理合起来正表明，连续映射 $f: x \to f(x)$ 把闭区间 $[a, b]$ 映成闭区间 $[m, M]$，其中 $m = \min\limits_{a \leqslant x \leqslant b} f(x)$，$M = \max\limits_{a \leqslant x \leqslant b} f(x)$．

\begin{xiti}
\item 设 $f(x)$ 在 $[a, b]$ 上连续，并且最大值点 $x_0$ 是唯一的，又设 $x_n \in [a, b]$，使 $\lim\limits_{n \to \infty} f(x_n) = f(x_0)$，求证
\[
\lim_{n \to \infty} x_n = x_0.
\]
\item 设 $f(x)$ 在 $[a, b]$ 上连续，可微；又设
\begin{ziti}
\item $\min\limits_{a \leqslant x \leqslant b} f(x) < p < \max\limits_{a \leqslant x \leqslant b} f(x)$；
\item 如果 $f(x) = p$，则有 $f'(x) \neq 0$，
\end{ziti}
求证：$f(x) = p$ 的根只有有限多个．
\item 设 $f(x)$ 在 $[a, b]$ 连续，$f(a) < 0$，$f(b) > 0$，求证：存在 $\xi \in (a, b)$，使 $f(\xi) = 0$，且 $f(x) > 0$（$\xi < x \leqslant b$）．
\item 设 $f(x)$ 是 $[a, b]$ 上的连续函数，其最大值和最小值分别为 $M$ 和 $m$（$m < M$），求证：必存在区间 $[\alpha, \beta]$，满足条件：
\begin{ziti}
\item $f(\alpha) = M$，$f(\beta) = m$ 或 $f(\alpha) = m$，$f(\beta) = M$；
\item $m < f(x) < M$，当 $x \in (\alpha, \beta)$．
\end{ziti}
\item 设 $f(x)$ 在 $[0, 2a]$ 连续，且 $f(0) = f(2a)$，求证：存在 $x \in [0, a]$，使 $f(x) = f(x + a)$．
\item 设 $f(x)$ 在 $[a, b]$ 上连续，且取值为整数，求证：$f(x) \equiv$ 常数．
\item 设 $f(x)$ 在 $(a, b)$ 上一致连续，$a, b \neq \pm\infty$，证明 $f(x)$ 在 $(a, b)$ 上有界．
\item 若函数 $f(x)$ 在 $(a, b)$ 上满足利普希茨（Lipschitz）条件，即存在常数 $K$，使得
\[
|f(x') - f(x'')| \leqslant K |x' - x''|, \quad x', x'' \in (a, b).
\]
证明：$f(x)$ 在 $(a, b)$ 上一致连续．
\item 试用一致连续的定义证明：若函数 $f(x)$ 在 $[a, c]$ 和 $[c, b]$ 上都一致连续，则 $f(x)$ 在 $[a, b]$ 上也一致连续．
\item 设 $f(x)$ 在 $(-\infty, +\infty)$ 上连续，且 $\lim\limits_{x \to -\infty} f(x)$ 与 $\lim\limits_{x \to +\infty} f(x)$ 存在，证明：$f(x)$ 在 $(-\infty, +\infty)$ 上一致连续．
\item 若 $f(x)$ 在区间 $X$（有穷或无穷）中具有有界的导数，即 $|f'(x)| \leqslant M$，$x \in X$，则 $f(x)$ 在 $X$ 中一致连续．
\item 求证：$f(x) = \sqrt{x} \ln x$ 在 $(0, +\infty)$ 上一致连续．
\item 设 $f(x)$ 在 $(a, +\infty)$ 上可导，且 $\lim\limits_{x \to +\infty} f'(x) = +\infty$，求证：$f(x)$ 在 $(a, +\infty)$ 上不一致连续．
\item 求证：$f(x) = x \ln x$ 在 $(0, +\infty)$ 上不一致连续．
\end{xiti}
